\section{Introduction}
\label{sec:introduction}

Recent advances in Large Language Models (LLMs), Vision-Language Models
(VLMs), and vision-action foundation models have enabled increasingly capable
agents for open-world games~\citep{wang2024voyager,magne2026nitrogen}, web
navigation~\citep{zhou2023webarena,deng2023mind2web}, and general computer
control~\citep{liu2024agentbench,tan2024cradle}. Unlike single-step question
answering, these environments require agents to act over long, changing
interaction histories with delayed feedback. In tasks such as
\stardojo{}~\citep{tan2025stardojo}, which spans 1,000 tasks and a Lite-100
subset, an agent may need dozens or hundreds of steps to coordinate tool use,
navigate, handle menus, switch subgoals, and recover from failures without a
privileged game-state API. The central challenge is therefore \textbf{closed-loop
control under scarce deliberation}: reasoning should be spent on bottlenecks and
recovery, while routine low-level actions must stay cheap. This creates a
\textbf{planning-latency-memory trilemma}: deeper reasoning improves
deliberation but raises inference cost, lightweight reactive execution is fast
but short-sighted, and naive memory reuse introduces retrieval noise.

Current agents address only parts of this control problem. As illustrated in
Figure~\mbox{\ref{fig:motivation}}-(a), planning-intensive frameworks such as
CRADLE can reason effectively, but repeated calls to large models make
high-frequency control expensive in tokens and latency~\citep{tan2024cradle}.
Conversely, Figure~\mbox{\ref{fig:motivation}}-(b) captures reactive agents such
as ReAct: they execute cheaply but become unreliable under drift, execution
failure, repetition, and multi-stage transitions~\citep{yao2023react}.
The desired behavior, shown in Figure~\mbox{\ref{fig:motivation}}-(c), reuses
strategic guidance across routine reactive steps while retaining the ability to
replan when the local state changes. Fast-slow and hierarchical agents, including
SayCan, Inner Monologue, and SwiftSage~\citep{ahn2022saycan,huang2022innermonologue,lin2023swiftsage},
as well as ADaPT and Optimus-3~\citep{prasad2024adapt,li2025optimus3}, move
toward selective reasoning. However, multimodal open worlds leave three
gaps: triggering is often rigid, sudden visual changes are hard to detect from
text-only cues, and the fast layer is frequently treated as a passive executor.
Meanwhile, long-horizon memory is often a flat prompt buffer or single
retrieval pool, weakening cross-episode reuse and increasing retrieval
noise~\mbox{\citep{park2023generativeagents,packer2023memgpt}}.

Motivated by these observations, we propose \method{}, an \textbf{event-triggered
deliberation framework} for cost-efficient long-horizon game agents. Our core
insight is that strategic reasoning should be amortized across locally stable
segments and re-invoked only at event boundaries. Long-horizon game control is
temporally uneven: most steps lie in stable local execution, while a few
boundaries dominate recoverable failures. This turns reasoning into a resource
allocation problem. \method{} therefore couples three mechanisms: the Event
Trigger decides \emph{when} expensive reasoning is worth invoking; the bounded
Reactive Controller decides \emph{how} to act safely between deliberations rather
than serving as a passive executor; and controller-specific Hierarchical Memory
decides \emph{what evidence} supports local action reuse versus strategic
replanning. As detailed in Section~\ref{sec:method}, the memory separates the
State-Action Memory Bank (\samb{}) from the State-Action Knowledge Graph
(SA-KG), preventing routine execution and replanning from sharing one noisy
retrieval context.

\begin{figure}[t!]
    \centering
    \includegraphics[width=0.98\textwidth]{taser.pdf}
    \caption{\textbf{Motivation for event-triggered deliberation.}
    Single-Strategic \textbf{\textit{(a)}} control reasons carefully but is expensive, while Single-
    Reactive \textbf{\textit{(b)}} control is cheap but brittle under drift and bottlenecks. \textbf{\textit{(c)}} \method{}
    allocates strategic reasoning only at likely bottlenecks, aiming for higher
    success and recovery with lower per-step cost.}
    \label{fig:motivation}
    \vspace{-0.6em}
\end{figure}

Overall, our main contributions are:
\begin{itemize}
    \item \textbf{Event-triggered amortized deliberation.} We formulate
    long-horizon multimodal control as budgeted allocation of strategic reasoning
    over event boundaries.
    \item \textbf{Adaptive dual-controller execution.} \method{} uses strategic
    planning at escalation points and bounded reactive override during stable
    local execution.
    \item \textbf{Controller-specific hierarchical memory.} \method{} separates
    \samb{} for local state-action reuse from SA-KG for structured strategic
    evidence.
    \item \textbf{Cost-aware evidence.} On \stardojo{} Lite-100 and RDR2,
    \method{} improves success--cost trade-offs, with ablations isolating trigger
    timing, override, and memory structure.
\end{itemize}
